% GENERATED, edit the French unit or the translation table
\mxnewpage\section*{EV2-B -- Mock exam 1}\label{EV2-B}
\addcontentsline{toc}{section}{EV2-B -- Mock exam 1}
\textbf{Paper reserved for the teacher}. Chapters 1 to 7.\par Duration: 2 h. No calculator. Marked out of 20. 14--16 January 2027.
\subsection*{Part 1 -- Fluency (6 points)}
Twelve questions, 0.5 point each. Only one correct answer; no justification required. An incorrect or missing answer scores 0, with no loss of points.
\begin{enumerate}
\item Calculate $5+2\times3$.\par A. $11$\quad B. $10$\quad C. $14$\quad D. $21$
\item Calculate $\frac12+\frac13$.\par A. $\frac{5}{3}$\quad B. $\frac{1}{6}$\quad C. $\frac{5}{6}$\quad D. $\frac{2}{5}$
\item Calculate $\dfrac{2^5}{2^{3}}$.\par A. $2$\quad B. $10$\quad C. $1$\quad D. $4$
\item A value of 50 increases by 20 \%. Final value?\par A. $10$\quad B. $60$\quad C. $40$\quad D. $70$
\item A value of 80 is obtained after a 20 \% decrease. Initial value?\par A. $100$\quad B. $64$\quad C. $96$\quad D. $60$
\item A 25 \% increase then a 20 \% decrease: overall coefficient?\par A. $\frac{9}{10}$\quad B. $1$\quad C. $\frac{21}{20}$\quad D. $\frac{4}{5}$
\item The solution of $2x+5=0$ is\par A. $-10$\quad B. $\frac{5}{2}$\quad C. $\frac{-5}{2}$\quad D. $10$
\item The slope between $(0;5)$ and $(2;11)$ is\par A. $2$\quad B. $-3$\quad C. $3$\quad D. $6$
\item The mean of $5, 7, 9$ is\par A. $7$\quad B. $21$\quad C. $8$\quad D. $5$
\item A sector of $90^\circ$ represents what fraction of the total?\par A. $\frac{1}{3}$\quad B. $\frac{1}{2}$\quad C. $\frac{3}{4}$\quad D. $\frac{1}{4}$
\item A probability is $\dfrac{2}{5}$. Its complement is\par A. $\frac{5}{2}$\quad B. $\frac{2}{5}$\quad C. $\frac{1}{5}$\quad D. $\frac{3}{5}$
\item For $f(x)=x^2+5$, $f(-2)$ is\par A. $1$\quad B. $7$\quad C. $9$\quad D. $20$
\end{enumerate}
\mxnewpage\subsection*{Part 2 -- Independent exercises (14 points)}
Every step must be justified. Exact results are expected unless stated otherwise.
\mxneed{12}
\subsection*{Exercise 1 -- Independence and repetition (5 points)}
Two events $A$ and $B$ are independent and satisfy $P(A)=\dfrac14$ and $P(B)=\dfrac23$. The first question concerns these events. In a separate experiment, used in questions 2 and 3, a trial with success probability $\dfrac14$ is repeated three times independently.
\begin{enumerate}
\item Calculate $P(A\cap B)$ and $P(A\cup B)$. \hfill(1 pt)
\item Calculate the probability of exactly two successes in the three trials. \hfill(2 pts)
\item Calculate that of at least one success and justify the move to the complement. \hfill(2 pts)
\end{enumerate}
\mxneed{12}
\subsection*{Exercise 2 -- Rational function and symmetry (5 points)}
The function $f$ is defined on $\R\setminus\{0\}$ by $f(x)=\dfrac{x^2+4}{x}$. Its domain is symmetric about 0, but is not an interval. We will use its parity to link the studies on the two half-lines.
\begin{enumerate}
\item Study the parity and rewrite f. \hfill(1 pt)
\item Study the variation on $]0;+\infty[$. \hfill(2 pts)
\item Deduce the variation on $]-\infty;0[$ and explain why f is not increasing on its whole domain. \hfill(2 pts)
\end{enumerate}
\mxneed{12}
\subsection*{Exercise 3 -- Distance and parameter (4 points)}
In an orthonormal frame, given the line $d$ with equation $3x+4y-12=0$ and the point $P(0;0)$. Let $H$ be the orthogonal projection of $P$ onto $d$. We look for the circle with centre $P$ tangent to this line.
\begin{enumerate}
\item Give a normal and check that $(4;0)$ is on d. \hfill(1 pt)
\item Calculate the projection H of P. \hfill(1 pt)
\item Determine the distance from P to d and the equation of the circle with centre P tangent to d. \hfill(2 pts)
\end{enumerate}
